\documentclass[11pt,a4paper]{article}

\usepackage[utf8]{inputenc}
\usepackage[T1]{fontenc}
\usepackage{geometry}
\geometry{left=2cm, right=2cm, top=2cm, bottom=2cm}
\usepackage{enumitem}
\usepackage{hyperref}
\usepackage{titlesec}
\usepackage{xcolor}

% Color setup for links
\hypersetup{
    colorlinks=true,
    linkcolor=blue,
    filecolor=magenta,      
    urlcolor=blue,
    pdftitle={Resume - Chenjie Mao},
}

% Section formatting
\titleformat{\section}{\large\bfseries\uppercase}{}{0em}{}[\titlerule]
\titlespacing{\section}{0pt}{1.5ex plus 1ex minus .2ex}{1ex plus .2ex}

\setlength{\parindent}{0pt}

\begin{document}

% Header
\begin{center}
    {\Huge \bfseries Chenjie Mao} \\ \vspace{0.5em}
    \href{mailto:cmao@wustl.edu}{cmao@wustl.edu}
\end{center}

\vspace{1em}

% About Me
\section*{About Me}
I am a PhD student at Washington University in St. Louis, advised by Professor Chongjie Zhang. My research interests encompass a broad range of topics in learning theory, including reinforcement learning, online learning, and deep learning theory. Recently, I have focused primarily on the statistical aspects of (offline) reinforcement learning, especially its non-asymptotic behaviors and extensions to reinforcement learning with human feedback (RLHF).

% Education
\section*{Education}
\textbf{Huazhong University of Science and Technology} \hfill 2020 -- 2024 \\
Bachelor of Engineering, School of Computer Science and Technology

\textbf{Huazhong University of Science and Technology} \hfill 2025 -- Present \\
Phd Student, Computer Science and Engineering

% Work Experience
\section*{Work Experience}
\textbf{Shanghai Artificial Intelligence Laboratory} \hfill 2023.7 -- 2025.6 \\
\textit{Research Assistant}
\begin{itemize}[leftmargin=*, noitemsep, topsep=2pt]
    \item Investigated the statistical aspects of offline reinforcement learning and related areas.
    \item Explored potential alignments between RL and RLHF.
\end{itemize}

% Preprints & Publications
\section*{Preprints \& Publications}
\begin{enumerate}[leftmargin=*, noitemsep, topsep=2pt]

    \item \textbf{Offline Reinforcement Learning with Additional Covering Distributions} \\
    Chenjie Mao\footnote{Completed independently as a third-year undergraduate student.} \\
    \textit{TMLR (Transactions on Machine Learning Research)}
    \begin{itemize}[leftmargin=1.5em, noitemsep, topsep=1pt]
        \item Proposed an algorithm with statistical guarantees for offline RL, utilizing general function approximation under single coverage and realizability-type assumptions, based on an additional covering distribution.
    \end{itemize}
    
    \vspace{0.5em}
    
    \item \textbf{On the Role of General Function Approximation in Offline Reinforcement Learning} \\
    Chenjie Mao, Qiaosheng Zhang, Zhen Wang, Xuelong Li \\
    \textit{ICLR 2024 Spotlight, Top 5\%}
    \begin{itemize}[leftmargin=1.5em, noitemsep, topsep=1pt]
        \item Investigated the statistical limitations of offline reinforcement learning with general function approximation and established some lower bounds.
    \end{itemize}

    \item \textbf{Rethinking the Hardness of PbRL: A Provable General Regret Bound} \\
    Chenjie Mao, Yi Fan, Ning Zhang, Chongjie Zhang\\
    \textit{International Conference on Machine Learning (ICML), 2026}
    \begin{itemize}[leftmargin=1.5em, noitemsep, topsep=1pt]
        \item Introduced RTPQ, a simple and provably efficient algorithm that achieves near-optimal $\sqrt{T}$ regret in preference-based reinforcement learning with trajectory-level feedback.
    \end{itemize}
\end{enumerate}

% Interests
\section*{Interests}
I have broad interests in topics in machine learning, such as:
\begin{itemize}[leftmargin=*, noitemsep, topsep=2pt]
    \item Reinforcement Learning
    \item Generative Models
\end{itemize}

% Skills
\section*{Skills}
\textbf{Proficient:} Python, PyTorch, Git, Linux, LaTeX \\
\textbf{Have Experience:} Shell Scripting, C/C++, Golang

\end{document}
